\documentclass[exercice,a4]{thedocument}%
\usepackage{texmaths-tikz}%
\usepackage{pstricks}%
\usepackage{pst-bezier}%
\usepackage{pstricks-add}%
\usepackage{pst-eucl}%
\usepackage{pst-func}%
\usepackage{pst-tree}%
\usepackage{pst-circ}%
\usepackage{tikz}%
\usepackage{pgfplots}%
\usepackage{tkz-tab}%
\usepackage{tikzpagenodes}%
\usepackage{tabularray}%
\usepackage{tabularx}%
\usepackage{multirow}%
\usepackage{multicol}%
\usepackage{scratch3}%
\usepackage{enumitem}%
\usepackage[np]{numprint}%
\usepackage{pas-tableur}%
\usepackage{listings}%
\usepackage{graphicx}%
\usepackage{adjustbox}%
\usepackage{needspace}%
\usepackage{xcolor}%
\usepackage{eurosym}%
\usepackage{diagbox}%
\usepackage{colortbl}%
\usepackage{multido}%
\usepackage{icomma}%
\usepackage[french]{babel}%
\graphicspath{{../brevet/2026/}}%
\startdatakeys
    \source{26GENMATMEAG3}
    \BO{2026}
    \cycle{4}
    \competencesDuSocle{CA,RA,CO}
    \objectifsApprentissage{C4A1aa,C4A2ad,C4B2cb,C4D2aa,C4D3ac}
    \motsCles{moyenne,pourcentage,Pythagore,trigonométrie,angle,vitesse,durée}
\enddatakeys
\definecolor{gris}{gray}{0.5}%
\providecommand{\wideparen}[1]{\overset{\frown}{#1}}%
\begin{document}%
\providecommand{\degres}{^\circ}%
\emph{Les parties {\rm A} et {\rm B} de cet exercice sont indépendantes.}

\medskip

Le tableau ci-dessous donne le nombre total de médailles obtenues par 3 nations européennes, l'Italie, la France et la Norvège, lors de six éditions des Jeux Olympiques d'hiver.

\begin{center}
\begin{tblr}{colspec={|X[c]|*{6}{X[c]|}},row{1}={font=\bfseries},column{1}={font=\bfseries}}
	\hline
	Médailles & 2002 & 2006 & 2010 & 2014 & 2018 & 2022 \\ \hline
	Italie    & 13   & 11   & 5    & 8    & 10   & 17   \\ \hline
	France    & 11   & 9    & 11   & 15   & 15   & 14   \\ \hline
	Norvège   & 25   & 19   & 23   & 27   & 39   & 37   \\ \hline
\end{tblr}
\end{center}

\subsection*{Partie A}

\begin{enumerate}
	\item Donner le nombre de médailles obtenues par la France en 2010.
	\item Calculer la moyenne des médailles obtenues par la France lors de ces 6 jeux olympiques.
	\item Lors des jeux olympiques de 2002, la Norvège a terminé à la première place des nations. Elle avait obtenu au total 25 médailles dont 13 médailles en or.

	Calculer le pourcentage de médailles d'or gagnées par la Norvège en 2002 par rapport au nombre total des médailles qu'elle a obtenues cette année-là.
\end{enumerate}

\subsection*{Partie B}

Pour accéder au point de départ d'une épreuve, les skieurs utilisent un téléphérique dont le schéma ci-dessous représente le déplacement entre les points A et B.

%Début Tikz
\begin{center}
\begin{tikzpicture}[line cap=round, line join=round]
	\coordinate (A) at (0,0);
	\coordinate (H) at (8,0);
	\coordinate (B) at (8,4.3);
	\draw[line width=1.2pt,red,fill=red!30] (H) rectangle ++(-0.4,0.4);
	\draw[line width=1.2pt] (A) -- (H) -- (B) -- cycle;
	\node[below left] at (A) {A};
	\node[below right] at (H) {H};
	\node[above] at (B) {B};
	\draw[<->, line width=0.8pt, > = stealth,shift={(118.26:0.5)}] (A) -- (B) node[pos=0.5,above,sloped]{\np{1600}~m};
\end{tikzpicture}
\end{center}
%Fin Tikz

La longueur AB est égale à \np{1600}~m.

L'altitude du point A est \np{1814}~m et celle du point B est \np{2359}~m.

Les points A et H sont à la même altitude.

\begin{enumerate}
	\item Montrer que la différence d'altitude entre les points H et B, c'est-à-dire la longueur HB, est égale à 545~m.
	\item Pour respecter les normes de sécurité du téléphérique, la mesure de l'angle $\widehat{\mathrm{BAH}}$ doit être inférieure à 30\degres.

	Cette norme est-elle respectée ?\\
	\textbf{Détailler la réponse en précisant la démarche.}
	\item La vitesse du téléphérique est égale à 5~m/s.

	Quelle est la durée du trajet entre A et B ? Donner le résultat en minutes et secondes.
\end{enumerate}

\needspace{10\baselineskip}
\medskip
\showthesource
\end{document}
